\section{Introduction}

Engineering and technology programmes increasingly emphasise the need to align student learning with professional practice. Across Europe and beyond, universities are encouraged to adopt active learning pedagogies, strengthen links with industry, and design learning environments that support the development of transferable professional competences alongside disciplinary knowledge.

This paper presents a case describing the redesign of an undergraduate Business Intelligence course. The redesign was informed by the author’s transition back into higher education after an extended period in industry and draws explicitly on professional practice as a source of course structure, assessment logic, and learning activities. The aim was to make the course both pedagogically coherent and professionally relevant, while remaining feasible within the organisational constraints of university teaching.

\section{Context and Challenges}

Business Intelligence is an interdisciplinary field that combines data analysis, information systems, and business decision-making. Graduates are expected to integrate technical competence with domain understanding, communicate results clearly, and collaborate effectively in teams. While demand for these competences is high, academic courses often focus on tools and techniques in isolation, with limited attention to professional working practices.

When the author first taught the course as a part-time instructor, it already incorporated elements of Team-Based Learning, Problem-Based Learning, and flipped classroom pedagogy. However, these approaches were only loosely connected, and the project component offered limited opportunity to engage with realistic industrial problems. Student feedback indicated high workload and weak coherence between course components, while industry collaborators reported that project outcomes had limited practical value. These observations pointed to a misalignment between learning activities, assessment, and professional expectations.

\section{Pedagogical Approach}

The redesigned course integrates flipped classroom pedagogy with Team-Based Learning (TBL) and Problem-Based Learning (PBL). In a flipped classroom approach, students engage with core material prior to class, allowing scheduled contact time to be used for discussion, problem solving, and application \cite{tucker2012flipped}. TBL structures this process through individual preparation, team readiness activities, and team-based application tasks \cite{michaelsen2004team,fink2003creating,sibley2014getting}.

\begin{figure}[ht]
\centering
\tikzset{
  compbox/.style={
    draw, rounded corners=5pt, align=center,
    text width=30mm, minimum height=8mm,
    font=\scriptsize, inner sep=1pt
  },
  highlight/.style={
    compbox,
    fill=green!20, draw=green!60!black, very thick
  }
}

\begin{tikzpicture}
  \matrix[matrix of nodes,
          nodes={compbox},
          column sep=2mm, row sep=2mm] {
    {Civic engagement} &
    {Creative thinking} &
    {Critical thinking} &
    {Ethical reasoning} \\
    {Global learning} &
    |[highlight]| {Information literacy} &
    {Inquiry and analysis} &
    {Integrative learning} \\
    {Intercultural knowledge \& competence} &
    {Foundations for life-long learning} &
    {Oral communication} &
    |[highlight]| {Problem solving} \\
    {Quantitative literacy} &
    {Reading} &
    |[highlight]| {Teamwork} &
    {Written communication} \\
  };
\end{tikzpicture}
\caption{The LOUIS competence framework, showing the three competences emphasised in the Business Intelligence course: problem solving, information literacy, and teamwork. See \cite{auroraLOUIS}.}
\label{fig:louis-grid}
\end{figure}

To support coherence between learning outcomes and assessment, the redesign was informed by the LOUIS competence framework developed within the Aurora university alliance \cite{auroraLOUIS}. From the framework, three competences were prioritised: problem solving, information literacy, and teamwork. These competences were embedded across learning activities and assessment and used as guidance when revising the course structure and project work. The relationship between the selected competences and the full LOUIS framework is shown in Fig.~\ref{fig:louis-grid}. The course design was also informed by the principles of Understanding by Design, which emphasise alignment between intended learning outcomes, teaching activities, and assessment \cite{wiggins2005understanding}.

\section{Industry Collaboration and Project Structure}

A central change in the redesign was the introduction of a single industrial context sustained throughout the semester. Instead of working with new datasets in each thematic block, students collaborated with one industrial partner and explored a shared dataset across all course activities. The partner contributed data, problem framing, and feedback on student deliverables, while students were encouraged to extend the dataset where relevant.

The course was organised into five thematic iterations, each focusing on a specific Business Intelligence method applied to the shared dataset. This allowed students to build familiarity with both the data and the problem context while concentrating on new analytical approaches in each iteration. The course concluded with a capstone project in which students integrated their prior work into a coherent solution for the industrial partner, consistent with established descriptions of capstone learning experiences \cite{capstone}. The overall structure of the course and the progression from thematic iterations to the final capstone project are summarised in Fig.~\ref{fig:course-flow}.


\begin{figure}[ht]
\centering
\begin{tikzpicture}
    [scale=0.62, transform shape,
    node distance=65mm,
    inner sep=3pt,
    auto,
    block/.style={
        draw,
        thick,
        fill=green!40,
        text width=30mm,
        align=center,
        minimum height=10mm
    },
    line/.style={
        draw, thick, ->, shorten >=2pt
    },
    dashedline/.style={
        draw, thick, ->, dashed, shorten >=2pt, color=black!50
    },
    week/.style={anchor=south west, text width=45mm},
    red/.style={block, draw=red, fill=red!40},
    orange/.style={block, draw=orange, fill=orange!40},
    yellow/.style={block, draw=yellow, fill=yellow!40},
    green/.style={block, draw=green, fill=green!40},
    violet/.style={block, draw=violet, fill=violet!40},
    blue/.style={block, draw=blue!40, fill=blue!20},
    agile/.style={block, draw=gray, fill=gray!10,inner sep=3mm}
    ]

    % Samantektarverkefni
    \node[week] (themecapstone) at (0,-75mm) {Week 11--14\\ Capstone Project};
    \node[blue, right of=themecapstone, xshift=70mm, label=below:\textbf{Project}] (capstone) {Team (tCAP)};
    % Endurgjöf milli teyma
    \draw [line] (capstone) edge[loop above] (capstone);
    \node[blue, right of=capstone, label=below:\textbf{Reflection}] (capref) {Student (iCAP)};
    \draw[line] (capstone) to (capref);

    % Litur/lotur
    \foreach \week/\theme/\color/\pos in {1--2/Data Ethics/red/1, 3--4 /Large Language Models/orange/2, 5--6/Supervised Learning/yellow/3, 7--8 /Clustering/green/4, 9--10/Process Mining/violet/5}{
        \node[week] (week\pos) at (0,-\pos*10mm) {Week \week\\\theme};
        \node[\color, right of=week\pos, xshift=-50+10*\pos,label=below:\textbf{Preparation}] (RAT\pos)  {Student (iRAT)\\Team (tRAT)};
        \node[\color, right of=RAT\pos,label=below:\textbf{Project}] (APP\pos) {Team (tAPP)};
        \node[\color, right of=APP\pos,label=below:\textbf{Reflection}] (REF\pos) {Student (iREF)};

        % Línur milli eininga
        \path [line] (RAT\pos) -- (APP\pos);
        \path [line] (APP\pos) -- (REF\pos);

        % Jafningjamat innan Teams
        \draw [line] (APP\pos) edge[loop above] (APP\pos);
        
        % Tengingar við samantektarverkefni
        \path [dashedline] (APP\pos.east) to[out=-20, in=70] (capstone.north east);
        \path [dashedline] (REF\pos.east) to[out=-30, in=40] (capstone.east);
    }
    \begin{scope}[on background layer]
    \node[agile,fit=(RAT1) (RAT5),label=above:{\textbf{Sprint planning}}] () {};
    \node[agile,fit=(APP1) (APP5),label=above:{\textbf{Sprint}}] () {};
    \node[agile,fit=(REF1) (REF5),label=above:{\textbf{Backlog}}] () {};
    \end{scope}
\end{tikzpicture}

\caption{Course structure and progression across five thematic iterations leading to the capstone project. Each iteration spans two weeks and focuses on a specific Business Intelligence method applied within a shared industrial context.}
\label{fig:course-flow}
\end{figure}

\section{Collaboration Tools and Peer Review}

To support collaboration and transparency, GitHub Classroom was introduced as the main project platform. The tool enabled version control, shared documentation, and peer review of code and reports, reflecting common professional practice. Student activity within the platform provided visible evidence of individual contribution and supported ongoing supervision.

Peer review was implemented using mechanisms inspired by industrial workflows, where code and documents are reviewed during development rather than only at submission. Within teams, students reviewed each other’s work and provided structured feedback. In later stages of the course, teams also reviewed the work of other teams. Assessment of peer review focused on the quality and constructiveness of the feedback rather than numerical scoring.

\section{Assessment and Feedback}

Assessment was designed to support iterative learning while maintaining clear quality standards. Deliverables from the thematic iterations were assessed using the same criteria as the final capstone project, but grading was scaled to emphasise formative feedback. This approach allowed students to understand expectations for final project quality without being penalised for work that was still under development.

Individual assessment was based on visible contribution in the collaboration platform and structured individual reflection. This combination was intended to address challenges related to fairness and uneven workload distribution commonly associated with team-based projects.

\section{Outcomes and Reflections}

Evidence from student evaluations, reflective writing, graduate feedback, and input from the industrial partner suggests increased student engagement and improved coherence across learning activities. Attendance and participation were consistently high, and students reported that working with realistic data and tools helped them understand how Business Intelligence methods are applied in professional contexts.

At the same time, collaboration with industry revealed differences in expectations between academic and professional settings. While academic assessment can tolerate minor errors in developmental work, industry partners placed greater emphasis on reliability and trustworthiness. Recognising and discussing these differences became an important part of students’ learning and professional preparation.

\section{Conclusion}

The redesign of the Business Intelligence course demonstrates how sustained industry collaboration, combined with active learning and competence-oriented assessment, can strengthen the connection between academic learning and professional practice. By aligning learning activities, assessment, and tools with realistic working practices, the course supported both technical learning and the development of transferable competences. The experience highlights the importance of making professional expectations explicit and of designing assessment and feedback mechanisms that support iterative learning while maintaining clear standards.